\chapter*{Conclusion générale}
\addcontentsline{toc}{chapter}{Conclusion générale}
\argossetpart{CONCLUSION GÉNÉRALE}
\argossetchapter{Conclusion générale}

Le transport maritime conteneurisé constitue l'une des composantes essentielles du commerce mondial\cite{unctad2025}. Pourtant, un conteneur sec, dépourvu d'alimentation et de connectivité embarquées, ne permet pas de suivre en continu l'intégrité physique de la marchandise durant le transit. Cette opacité, source de litiges et de pertes financières, se double d'une contrainte réglementaire renforcée par les amendements SOLAS/MARPOL de janvier 2026\cite{omi2026}, particulièrement sensible sur les corridors camerounais et centrafricains\cite{corridor_transit,logcluster}. C'est dans ce contexte que s'inscrit le projet ARGOS, boîte noire logistique connectée visant à assurer la continuité numérique et l'intégrité physique de la chaîne conteneurisée, de la mer jusqu'à la destination finale.

Ce mémoire a conduit une démarche en trois volets. L'étude de marché a validé la pertinence commerciale du projet en substituant à la vente de matériel un modèle de service indexé sur la valeur de la marchandise (3\,\%), levant la barrière à l'entrée pour les transporteurs et instaurant une relation tripartite entre client finaux, propriétaires de conteneurs et promoteurs d'ARGOS. L'étude de faisabilité financière a confirmé la robustesse économique du projet, avec une VAN positive (Tableau~\ref{tab:van}), un TRI de 89,5\,\% et un délai de récupération d'environ vingt mois (Tableau~\ref{tab:tri_drci}). Enfin, l'étude technique et la réalisation du prototype ont démontré la faisabilité opérationnelle : architecture logicielle complète, base de données de 48 tables et 559 champs, chaîne d'acquisition embarquée (température, humidité, choc, position GPS, audio) reliée à une plateforme web de supervision, évaluée par les essais du Tableau~\ref{tab-validation-prototype}.

L'apport principal de ce travail réside dans le caractère intégré de la solution : ARGOS associe traçabilité, intégrité physique, preuve numérique horodatée et modèle économique adapté aux réalités d'Afrique centrale, où l'autonomie énergétique et la couverture réseau intermittente sont des contraintes majeures. Le projet répond ainsi aux trois objectifs fixés en introduction : sécurisation du conteneur en temps réel, réduction des pertes liées aux litiges et prise en compte des instruments internationaux pertinents, notamment SOLAS et le code IMDG\cite{imo_solas,imo_imdg}.

Ce travail comporte des limites assumées, qui tracent les développements futurs. L'intégration matérielle et les essais réels des modems LTE-M/NB-IoT et LoRaWAN restent à mener au-delà de la simulation ; la transmission dépend aujourd'hui des réseaux cellulaires terrestres, d'où des zones d'ombre qu'un module hybride cellulaire-satellitaire viendrait combler. L'infrastructure logicielle repose sur des services d'hébergement tiers ; pour garantir la souveraineté numérique des données du fret, un serveur physique local dédié devrait être déployé. Le service d'enrichissement BIC écrit encore des noms de colonnes obsolètes (\textit{type\_code}, \textit{tare\_mass\_kg}, \textit{max\_gross\_mass\_kg}) alors que le schéma actif expose \textit{iso\_type\_code}, \textit{tare\_kg} et \textit{max\_gross\_kg} ; ce contrat de données doit être aligné avant de considérer l'enrichissement comme fiable. L'interface, limitée au français et à l'anglais, gagnerait à être internationalisée pour les équipages et opérateurs portuaires transfrontaliers. Bien que les données générées soient scellées et auditables, leur force probante auprès des assurances et autorités suppose encore audits et certifications ; l'OMI établit les conventions maritimes mais ne certifie pas directement un boîtier commercial\cite{imo_solas,imo_imdg}. Enfin, les fonctionnalités décisionnelles reposant sur des algorithmes génériques, des travaux visant à entraîner un modèle d'IA dédié aux réalités du secteur portuaire.

En définitive, le projet ARGOS démontre qu'une solution de traçabilité active, techniquement maîtrisée et économiquement viable, peut transformer le conteneur passif en un objet connecté et digne de confiance. Les prochaines étapes telque la validation en conditions réelles, l'homologations et la spécialisation d'un modèle d'IA dédié ouvrent la voie à un déploiement à plus grande échelle sur les corridors logistiques d'Afrique centrale, et à une contribution africaine aux enjeux mondiaux de sécurisation du fret maritime.
